% GENERATED FILE. Edit canonical lessons, then run npm run generate:books.
% canonical-source-hash: fnv1a64:7b2ad8fc12eb4607
% canonical-lessons: RU-C91-west, RU-C91-underground, RU-C91-tram, RU-C91-ship, RU-C91-boat

\chapter{West, Underground, Tram, Ship, Boat}
\label{ch:ru-west-underground-tram-ship-boat}
\hlchaptermodality{91}

\begin{chapteropening}
\textbf{By the end of this chapter:} \emph{I can say the west, the underground, a tram, a ship and a boat.}

Complete the last lesson of chapter 91: I can say the west, the underground, a tram, a ship and a boat.
\end{chapteropening}

\section[\texorpdfstring{\ru{запад} (západ)}{запад (západ)}]{\ru{запад} (západ) — the west}

\label{lesson:RU-C91-west}



\begin{quote}
Before the new one: say the Russian for the south, then the Russian for the east.
\end{quote}

\subsection*{You'll want to know: \ru{запад}}
\textbf{\ru{запад}} (\emph{západ}) — \textquotedblleft{}the west\textquotedblright{}.

\begin{cousinweb}
\textbf{\ru{запад}} is \emph{\ru{за}-}, \textquotedblleft{}behind\textquotedblright{}, with \emph{\ru{пад}}, from \emph{\ru{падать}}, \textquotedblleft{}to fall\textquotedblright{}: first \textquotedblleft{}sunset\textquotedblright{}, like Latin \emph{occidēns}, \textquotedblleft{}the falling (sun)\textquotedblright{} (Vasmer).
\end{cousinweb}

\begin{grammarlens}[title={What you've built}]
One more word for this chapter.
\end{grammarlens}

\subsection*{Your turn}
\begin{itemize}
\raggedright
  \item \emph{Say it:} \emph{západ}
  \item \emph{Say it:} \emph{západ}, once more
  \item \emph{Say it:} say \emph{západ}
  \item \emph{Recall:} say the Russian for the south, then the Russian for the east, then say \emph{západ} again
\end{itemize}

\begin{tcolorbox}[breakable,colback=teal!4,colframe=teal!35!black,title={Before you move on}]
What does \ru{запад} mean? (\textquotedblleft{}The west\textquotedblright{}.) Say it once more.
\end{tcolorbox}
\section[\texorpdfstring{\ru{метро} (metró)}{метро (metró)}]{\ru{метро} (metró) — the underground}

\label{lesson:RU-C91-underground}



\begin{quote}
Before the new one: say the Russian for the east, then the Russian for the west.
\end{quote}

\subsection*{You'll want to know: \ru{метро}}
\textbf{\ru{метро}} (\emph{metró}) — \textquotedblleft{}the underground\textquotedblright{}.

\begin{cousinweb}
\textbf{\ru{метро}} comes from French \emph{métro}, cut short from \emph{métropolitain}, \textquotedblleft{}of the capital\textquotedblright{} (Vasmer).
\end{cousinweb}

\begin{grammarlens}[title={What you've built}]
One more word for this chapter.
\end{grammarlens}

\subsection*{Your turn}
\begin{itemize}
\raggedright
  \item \emph{Say it:} \emph{metró}
  \item \emph{Say it:} \emph{metró}, once more
  \item \emph{Say it:} say \emph{metró}
  \item \emph{Recall:} say the Russian for the east, then the Russian for the west, then say \emph{metró} again
\end{itemize}

\begin{tcolorbox}[breakable,colback=teal!4,colframe=teal!35!black,title={Before you move on}]
What does \ru{метро} mean? (\textquotedblleft{}The underground\textquotedblright{}.) Say it once more.
\end{tcolorbox}
\section[\texorpdfstring{\ru{трамвай} (tramváy)}{трамвай (tramváy)}]{\ru{трамвай} (tramváy) — a tram}

\label{lesson:RU-C91-tram}



\begin{quote}
Before the new one: say the Russian for the west, then the Russian for the underground.
\end{quote}

\subsection*{You'll want to know: \ru{трамвай}}
\textbf{\ru{трамвай}} (\emph{tramváy}) — \textquotedblleft{}a tram\textquotedblright{}.

\begin{cousinweb}
\textbf{\ru{трамвай}} came through French \emph{tramway} or straight from English \emph{tramway}: \emph{tram}, \textquotedblleft{}a beam, a rail\textquotedblright{}, with \emph{way} (Vasmer).
\end{cousinweb}

\begin{grammarlens}[title={What you've built}]
One more word for this chapter.
\end{grammarlens}

\subsection*{Your turn}
\begin{itemize}
\raggedright
  \item \emph{Say it:} \emph{tramváy}
  \item \emph{Say it:} \emph{tramváy}, once more
  \item \emph{Say it:} say \emph{tramváy}
  \item \emph{Recall:} say the Russian for the west, then the Russian for the underground, then say \emph{tramváy} again
\end{itemize}

\begin{tcolorbox}[breakable,colback=teal!4,colframe=teal!35!black,title={Before you move on}]
What does \ru{трамвай} mean? (\textquotedblleft{}A tram\textquotedblright{}.) Say it once more.
\end{tcolorbox}
\section[\texorpdfstring{\ru{корабль} (korábl')}{корабль (korábl')}]{\ru{корабль} (korábl') — a ship}

\label{lesson:RU-C91-ship}



\begin{quote}
Before the new one: say the Russian for the underground, then the Russian for a tram.
\end{quote}

\subsection*{You'll want to know: \ru{корабль}}
\textbf{\ru{корабль}} (\emph{korábl'}) — \textquotedblleft{}a ship\textquotedblright{}.

\begin{cousinweb}
\textbf{\ru{корабль}} — an old loan from Greek \emph{karábion}, \emph{kárabos}, \textquotedblleft{}a boat\textquotedblright{}, first \textquotedblleft{}a crab\textquotedblright{}; Vasmer dates it early because it has \emph{\ru{б}}, not \emph{\ru{в}}. Spanish \emph{carabela}, \textquotedblleft{}caravel\textquotedblright{}, has the same source.
\end{cousinweb}

\begin{grammarlens}[title={What you've built}]
One more word for this chapter.
\end{grammarlens}

\subsection*{Your turn}
\begin{itemize}
\raggedright
  \item \emph{Say it:} \emph{korábl'}
  \item \emph{Say it:} \emph{korábl'}, once more
  \item \emph{Say it:} say \emph{korábl'}
  \item \emph{Recall:} say the Russian for the underground, then the Russian for a tram, then say \emph{korábl'} again
\end{itemize}

\begin{tcolorbox}[breakable,colback=teal!4,colframe=teal!35!black,title={Before you move on}]
What does \ru{корабль} mean? (\textquotedblleft{}A ship\textquotedblright{}.) Say it once more.
\end{tcolorbox}
\section[\texorpdfstring{\ru{лодка} (lódka)}{лодка (lódka)}]{\ru{лодка} (lódka) — a boat}

\label{lesson:RU-C91-boat}



\begin{quote}
Before the new one: say the Russian for a tram, then the Russian for a ship.
\end{quote}

\subsection*{You'll want to know: \ru{лодка}}
\textbf{\ru{лодка}} (\emph{lódka}) — \textquotedblleft{}a boat\textquotedblright{}.

\begin{cousinweb}
\textbf{\ru{лодка}} is old \emph{lodŭka}, a diminutive of \emph{\ru{лодья}}, now \emph{\ru{ладья}}, \textquotedblleft{}a boat\textquotedblright{}, from Proto-Slavic \emph{*oldĭji} — inherited. Vasmer relates it to Lithuanian \emph{aldijà}, \textquotedblleft{}canoe\textquotedblright{}, and Swedish \emph{ålla}, \textquotedblleft{}a trough\textquotedblright{}.
\end{cousinweb}

\begin{grammarlens}[title={What you've built}]
One more word for this chapter.
\end{grammarlens}

\subsection*{Your turn}
\begin{itemize}
\raggedright
  \item \emph{Say it:} \emph{lódka}
  \item \emph{Say it:} \emph{lódka}, once more
  \item \emph{Say it:} say \emph{lódka}
  \item \emph{Recall:} say the Russian for a tram, then the Russian for a ship, then say \emph{lódka} again
\end{itemize}

\begin{tcolorbox}[breakable,colback=teal!4,colframe=teal!35!black,title={Before you move on}]
What does \ru{лодка} mean? (\textquotedblleft{}A boat\textquotedblright{}.) Say it once more.
\end{tcolorbox}
